\chapter{Rééquilibrage des indices}\label{ch:s1:index-rebalancing}

Lorsqu'une action entre dans un indice, chaque fonds qui suit l'indice doit l'acheter à la même clôture. Shleifer a montré en 1986 que les actions nouvellement ajoutées au S\&P 500
ont enregistré une rentabilité anormale significative à l'annonce, liée à leurs achats~: les courbes de demande des actions sont décroissantes. Trente ans plus tard
Patel et Welch ont constaté que l'effet s'inverse dans les années 2000~: les actions ajoutées et retirées ne subissent plus de déplacement permanent de la demande. Sur le marché synthétique, les \omterm{def:s1:index-rebalancing:fund}{fonds indiciels}
qui suivent 15~\% des actions de chaque membre poussent le cours d'un ajout de 5.3~\% en moyenne~; un trader qui achète à l'annonce en conserve 2.4~\% lorsque les prédicteurs
ont déjà pris la moitié du mouvement avant l'annonce, et 0.8~\% lorsqu'ils en ont pris les quatre cinquièmes. Ce chapitre porte sur les règles des indices, sur la prévision de ce
qu'elles feront et sur le trading d'un événement que tout le monde voit venir. La construction est \texttt{firm.indexevent}.

\section{Règles des indices et calendrier}

\begin{definition}[Fonds indiciel, jour de classement]\label{def:s1:index-rebalancing:fund}
Un \emph{fonds indiciel}\index{fonds indiciel} (index fund) détient les titres d'un indice dans les pondérations de celui-ci, de sorte que ses transactions suivent les changements de l'indice plutôt que
toute opinion sur les titres. Un \emph{jour de classement}\index{jour de classement} (rank day) est la date à laquelle un indice fondé sur des règles mesure les caractéristiques (capitalisation, flottant,
liquidité) qui déterminent l'appartenance à l'indice lors de la reconstitution suivante.
\end{definition}

Les changements d'indices (livre~1, chapitre~15) se répartissent en deux catégories. Les indices dont les changements sont décidés au cas par cas, comme le S\&P 500 dans les études ci-dessous, annoncent chaque changement avant qu'il ne prenne
effet.
Les indices fondés sur des règles, comme les indices Russell US, suivent un calendrier publié, et quiconque dispose des données peut en calculer les changements.

\begin{dated}{2026-09}{Le calendrier de reconstitution Russell}\label{dat:s1:index-rebalancing:russell}
FTSE Russell a rétabli une reconstitution semestrielle des indices Russell US en 2026. La reconstitution de juin 2026 a retenu le 30 avril comme \omterm{def:s1:index-rebalancing:fund}{jour de classement},
a publié les listes préliminaires le 22 mai et a pris effet après la clôture du vendredi 26 juin. La reconstitution de décembre 2026 a son \omterm{def:s1:index-rebalancing:fund}{jour de classement} le
vendredi 30 octobre, les ajouts et retraits préliminaires après la clôture du 13 novembre, un gel à partir du 30 novembre, et prend effet après la clôture du vendredi
11 décembre. FTSE Russell évalue à environ 12.2~billions de dollars les actifs d'investisseurs référencés sur les indices Russell US ou investis dans des produits qui les répliquent.
\end{dated}

L'indice synthétique suit le même schéma. Il regroupe les 300 plus grandes actions du marché synthétique, est reconstitué tous les 126 jours de négociation (dix-sept fois
sur l'échantillon), et encadre l'appartenance~: un non-membre n'entre que s'il se classe au 285e rang ou mieux au \omterm{def:s1:index-rebalancing:fund}{jour de classement}, et un membre ne sort que s'il se classe au-delà du 315e rang.
La bande (la règle de tampon du livre~1, chapitre~15) évite que les actions proches du seuil entrent et sortent sans cesse. Les listes préliminaires paraissent seize
jours de négociation après le \omterm{def:s1:index-rebalancing:fund}{jour de classement}, et les changements prennent effet à la clôture quarante jours de négociation après celui-ci~: 331 ajouts et 270 retraits au total.

\section{Prévision des ajouts et des retraits}

\begin{definition}[Prévision des entrées et sorties d'indice]\label{def:s1:index-rebalancing:prediction}
La \emph{prévision des entrées et sorties d'indice}\index{prévision des entrées et sorties d'indice} (membership prediction) estime, avant le \omterm{def:s1:index-rebalancing:fund}{jour de classement}, la probabilité que chaque action figure dans l'indice après la
reconstitution, en appliquant les règles de l'indice aux caractéristiques de l'action projetées dans le temps.
\end{definition}

Le \omterm{def:s1:index-rebalancing:fund}{jour de classement} lui-même, un changement fondé sur des règles est connu exactement~: l'application de la règle aux capitalisations de ce jour reproduit chaque ajout et chaque retrait. Avant
ce jour-là, les capitalisations doivent encore évoluer. \texttt{firm.indexevent.predict} simule l'évolution de la capitalisation de chaque action avec sa propre volatilité, deux
cents fois, et applique la règle à chaque tirage~; une action est prévue comme entrante si elle entre dans l'indice dans plus de la moitié des tirages.

\begin{center}
\small
\begin{tabular}{@{}lcccc@{}}
\toprule
prévision avant le \omterm{def:s1:index-rebalancing:fund}{jour de classement} & ajouts~: précision & rappel & retraits~: précision & rappel\\
\midrule
au \omterm{def:s1:index-rebalancing:fund}{jour de classement} & 100~\% & 100~\% & 100~\% & 100~\%\\
20 jours de négociation avant & 86~\% & 72~\% & 86~\% & 74~\%\\
60 jours de négociation avant & 76~\% & 45~\% & 76~\% & 46~\%\\
\bottomrule
\end{tabular}
\end{center}

À vingt jours, 86~\% des ajouts prévus se réalisent et 72~\% des ajouts effectifs avaient été prévus~; à soixante jours, les trois quarts des prévisions
se réalisent encore, mais moins de la moitié des ajouts sont repérés. Le compromis est celui du chapitre~: plus la prévision est précoce, moins elle est sûre, et plus la
part du mouvement de cours qui reste à venir est grande.

\section{Trading de l'événement et du retournement}

L'événement synthétique est injecté dans la rentabilité ajustée du marché de chaque action. Les \omterm{def:s1:index-rebalancing:fund}{fonds indiciels} suivent une part des actions de chaque membre et négocient les changements à la
clôture effective~; leur effet de demande sur le cours relève de la loi de la racine carrée du livre~7, chapitre~27, $0.7\,\sigma\sqrt{Q/V}$, où $Q$ désigne le nombre d'actions qu'ils doivent acheter et
$V$ le volume quotidien de l'action. Une part de l'effet (\texttt{early}) est prise avant l'annonce par les traders qui ont prévu le changement~; le reste s'accumule
entre l'annonce et la clôture effective~; la moitié de ce mouvement s'inverse au cours des vingt jours suivants. Avec 15~\% des actions suivies, la hausse moyenne est de 5.3~\%.

\begin{omfigure}
\begin{tikzpicture}
\begin{axis}[width=11cm,height=5.4cm,xlabel={\small jours de négociation depuis l'annonce},ylabel={\small rentabilité anormale cumulée (\%)},
  tick label style={font=\footnotesize},xmin=-20,xmax=44,ymin=-0.5,ymax=7,grid=major,grid style={gray!25},legend pos=south east,
  legend style={font=\scriptsize},table/col sep=comma]
\addplot[omThm,thick] table[x=day,y=early02]{figdata/strategies-1/10-index-rebalancing/path.csv};
\addplot[omDef,very thick] table[x=day,y=early05]{figdata/strategies-1/10-index-rebalancing/path.csv};
\addplot[black!50,thick,dashed] table[x=day,y=early08]{figdata/strategies-1/10-index-rebalancing/path.csv};
\addplot[black!40,dotted] coordinates {(0,-0.5) (0,7)};
\addplot[black!40,dotted] coordinates {(24,-0.5) (24,7)};
\legend{part anticipée 0.2,part anticipée 0.5,part anticipée 0.8}
\end{axis}
\end{tikzpicture}

\omcaption{Ajouts à l'indice synthétique, 15~\% des actions suivies~: rentabilité moyenne ajustée du marché avec l'effet des \omterm{def:s1:index-rebalancing:fund}{fonds indiciels} injecté, depuis vingt jours de négociation
avant l'annonce (jour 0) jusqu'à la clôture effective (jour 24) et les vingt jours suivants, pour trois parts de l'effet prises avant l'annonce.
Données~: \texttt{s1\_index.path\_for\_figure}.}\label{fig:s1:index-rebalancing:path}
\end{omfigure}

Trois stratégies se partagent l'événement. \emph{Prévoir} achète les ajouts (et vend les retraits) vingt jours avant l'annonce et tient la position jusqu'au bout~: avec une prévision parfaite
elle rapporte 3.75~\% par événement avec une part anticipée de moitié. \emph{Trading de l'annonce} achète à la clôture de l'annonce et vend à la clôture
effective aux \omterm{def:s1:index-rebalancing:fund}{fonds indiciels}~: 2.40~\% par événement après coûts (erreur type 0.4~\%). \emph{Fournir de la liquidité à la clôture effective}, en vendant aux \omterm{def:s1:index-rebalancing:fund}{fonds indiciels} et
en rachetant pendant les vingt jours suivants, gagne le retournement~: 2.65~\%. Les ajouts conservent une petite dérive propre, puisque les actions qui ont grossi jusqu'à entrer dans un indice
sont des gagnants récents, et le marché synthétique injecte du momentum~; Chen, Noronha et Singal ont observé l'asymétrie réelle dans le sens inverse, avec une hausse permanente
pour les ajouts et aucune baisse permanente pour les retraits, asymétrie qu'ils ont attribuée en partie à la sensibilisation des investisseurs.

\section{Comment la stratégie s'est érodée}

\begin{center}
\small
\begin{tabular}{@{}lccc@{}}
\toprule
rentabilité entre l'annonce et la clôture effective & 5~\% des actions suivies & 15~\% & 30~\%\\
\midrule
part anticipée 0.2 & 2.2~\% & 4.0~\% & 5.7~\%\\
part anticipée 0.5 & 1.3~\% & 2.4~\% & 3.5~\%\\
part anticipée 0.8 & 0.4~\% & 0.8~\% & 1.3~\%\\
\midrule
hausse moyenne & 3.1~\% & 5.3~\% & 7.5~\%\\
\bottomrule
\end{tabular}
\end{center}

Le rendement de l'événement croît avec la demande et décroît avec la part prise en amont. Davantage de capitaux indexés accroît la poussée~: de 3.1~\% pour 5~\% des actions suivies à
7.5~\% pour 30~\%. Plus de capital d'arbitrage qui anticipe le changement le fait remonter plus tôt, hors de portée de quiconque attend l'annonce~: à une part anticipée de 0.8
la transaction d'annonce garde un cinquième de ce qu'elle garde à 0.2. Le rendement propre des prédicteurs monte à 5.3~\% par événement, mais seulement avec une prévoyance parfaite~; vingt
jours à l'avance, un prédicteur se trompe sur une entrée sur sept et en manque plus d'un quart. C'est la configuration que montre le registre public pour le S\&P 500~: un
effet indice d'abord documenté comme un déplacement permanent de la demande, puis comme un déplacement temporaire qui se résorbe, sur un marché où l'événement est connu de tous. La
transaction n'a pas disparu~; elle s'est déplacée plus tôt, vers ceux qui prédisent, et plus tard, vers ceux qui fournissent de la liquidité à la clôture.

\section{Fiches de stratégie}

\begin{strategyfile}{Entrées prévues}\label{strat:s1:index-rebalancing:additions}
\sfield{Qui paie, et pourquoi} Les \omterm{def:s1:index-rebalancing:fund}{fonds indiciels} qui doivent acheter à la clôture effective, quel que soit le prix.
\sfield{Instruments et plateformes} Les actions proches du seuil de l'indice~; l'enchère de clôture à la date d'effet.
\sfield{Signal} La probabilité d'entrer, déduite des règles de l'indice appliquées aux caractéristiques projetées.
\sfield{Dimensionnement et exécution} Acheter les entrées probables avant le \omterm{def:s1:index-rebalancing:fund}{jour de classement} ou l'annonce~; vendre dans la demande des \omterm{def:s1:index-rebalancing:fund}{fonds indiciels} à la clôture effective.
\sfield{Coûts} Négociation sur des actions proches du seuil, souvent moins liquides~; prévisions erronées.
\sfield{Comment elle s'épuise} La concurrence qui fait avancer le cours~; \omterm{def:s1:index-rebalancing:fund}{fonds indiciels} qui négocient sur plusieurs jours.
\sfield{Horizon, capacité, infrastructure} Quelques semaines~; capacité limitée par la taille de l'événement~; un moteur de règles d'indice et des données de capitalisation.
\sfield{Backtest en toute rigueur} Règles telles qu'elles étaient à chaque date~; prévisions fondées sur des données antérieures au \omterm{def:s1:index-rebalancing:fund}{jour de classement}~; prévisions erronées comptées.
\sfield{Sources} Shleifer (1986)~; la simulation de ce chapitre (3.75~\% par événement avec une prévision parfaite à une part anticipée de la moitié).
\end{strategyfile}

\begin{strategyfile}{Sorties prévues}\label{strat:s1:index-rebalancing:deletions}
\sfield{Qui paie, et pourquoi} Les \omterm{def:s1:index-rebalancing:fund}{fonds indiciels} qui doivent vendre à la clôture effective.
\sfield{Instruments et plateformes} Les membres proches du seuil de sortie.
\sfield{Signal} La probabilité de sortir de l'indice.
\sfield{Dimensionnement et exécution} Vendre à découvert les sorties probables (si l'emprunt est possible) et les racheter aux \omterm{def:s1:index-rebalancing:fund}{fonds indiciels} qui vendent à la clôture.
\sfield{Coûts} Emprunt sur des actions en baisse, souvent de petite capitalisation.
\sfield{Comment elle s'épuise} Comme pour les entrées~; les actions retirées peuvent ne pas rester en baisse.
\sfield{Horizon, capacité, infrastructure} Quelques semaines~; disponibilité de l'emprunt.
\sfield{Backtest en toute rigueur} Frais d'emprunt et disponibilité~; radiation des membres les plus faibles.
\sfield{Sources} Chen, Noronha et Singal (2004)~: aucun recul permanent pour les sorties.
\end{strategyfile}

\begin{strategyfile}{Fourniture de liquidité à la clôture effective}\label{strat:s1:index-rebalancing:liquidity}
\sfield{Qui paie, et pourquoi} Les \omterm{def:s1:index-rebalancing:fund}{fonds indiciels} qui négocient à une seule clôture, quel que soit le prix.
\sfield{Instruments et plateformes} L'enchère de clôture (livre~1, chapitre~13) à la date d'effet.
\sfield{Signal} Le déséquilibre publié de l'enchère de clôture et la demande indicielle connue.
\sfield{Dimensionnement et exécution} Vendre les entrées et acheter les sorties à l'enchère~; procéder au débouclage sur les semaines suivantes.
\sfield{Coûts} Risque de détention pendant le retournement.
\sfield{Comment elle s'épuise} Lorsque les \omterm{def:s1:index-rebalancing:fund}{fonds indiciels} étalent leurs négociations, ou lorsque trop de capitaux se présentent à la clôture.
\sfield{Horizon, capacité, infrastructure} De quelques jours à quelques semaines~; données sur le déséquilibre de l'enchère.
\sfield{Backtest en toute rigueur} Cours de l'enchère, et non moyenne de la journée~; capacité rapportée à la taille de l'enchère.
\sfield{Sources} La simulation de ce chapitre (2.65~\% par événement lorsque la moitié de la poussée se retourne)~; Patel et Welch (2017) sur le retournement dans les années 2000.
\end{strategyfile}

\begin{strategyfile}{Retournement après l'événement}\label{strat:s1:index-rebalancing:reversal}
\sfield{Qui paie, et pourquoi} Comme pour la fourniture de liquidité~: la part de la poussée qui relevait de la pression, et non de l'information.
\sfield{Instruments et plateformes} Actions récemment entrées et sorties.
\sfield{Signal} La taille de la poussée rapportée au volume de l'action.
\sfield{Dimensionnement et exécution} Vendre à découvert les entrées et acheter les sorties après la date d'effet, avec couverture.
\sfield{Coûts} Emprunt sur les entrées~; le momentum des gagnants récents.
\sfield{Comment elle s'épuise} Si l'appartenance à l'indice véhicule de l'information (notoriété, liquidité), la poussée ne se retourne pas.
\sfield{Horizon, capacité, infrastructure} De quelques semaines à quelques mois.
\sfield{Backtest en toute rigueur} Séparer les entrées des sorties~; contrôler le momentum.
\sfield{Sources} Patel et Welch (2017)~; Chen, Noronha et Singal (2004).
\end{strategyfile}

\section{Tutoriel~: tout le monde le prévoit}\label{tut:s1:index-rebalancing:tut}

\begin{tutorial}
\textbf{Objectif.} Écrire une règle d'indice à bande, prévoir ses changements et négocier l'événement à trois moments. \textbf{État final~:} les deux tableaux et
la \cref{fig:s1:index-rebalancing:path}.
\begin{tutsteps}
\item \textbf{La règle et sa prévision}~: rangs, bande et simulation prospective des capitalisations.
\omcode{code/firm/indexevent/firm_indexevent.py}{27}{48}{La règle de reconstitution à bande et sa prévision.}\label{lst:s1:index-rebalancing:rule}
\item \textbf{L'événement}~: le chemin de cours simulé, de la prévision au retournement.
\omcode{code/firm/indexevent/firm_indexevent.py}{55}{63}{Le chemin de cours de l'événement.}\label{lst:s1:index-rebalancing:path}
\item \textbf{Lancer} \texttt{accuracy} à 0, 20 et 60 jours, \texttt{trades} sur la grille, et \texttt{fig\_index.py}.
\end{tutsteps}
\textbf{Pour aller plus loin.} Ne négocier que les prévisions au-delà d'une probabilité de 0.8 et compter les erronées~; laisser les \omterm{def:s1:index-rebalancing:fund}{fonds indiciels} négocier sur trois jours~; ajouter un
calendrier semestriel et un calendrier annuel, et comparer leur rotation.
\end{tutorial}

\section{Construction~: événements d'indice}\label{bld:s1:index-rebalancing:indexevent}

\begin{build}
\bfield{Rôle} Règles d'indice, leur prévision et le chemin de cours de la demande des \omterm{def:s1:index-rebalancing:fund}{fonds indiciels}, pour étudier les événements d'indice dont la vérité est connue.
\bfield{Interface} \texttt{reconstitute(cap, members, size, band)}, \texttt{predict(cap, vol, members, size, band, days, sims, rng)},
\texttt{event\_push(q, adv, sigma, eta)}, \texttt{event\_path(push, early, revert, pre, run, post)}.
\bfield{Règles} Appartenance déduite des seules données du \omterm{def:s1:index-rebalancing:fund}{jour de classement}~; prévisions fondées sur des données antérieures~; chemin de l'événement ajouté aux rendements ajustés du marché.
\bfield{Tests d'acceptation} \texttt{code/firm/indexevent/tests/}~: la règle de bande vérifiée à la main~; des prévisions certaines loin du seuil~; la poussée et le chemin vérifiés à la main.
\bfield{Pour aller plus loin} Ajustement au flottant et filtres de liquidité~; changements discrétionnaires de type comité~; \omterm{def:s1:index-rebalancing:fund}{fonds indiciels} qui négocient sur plusieurs jours.
\end{build}

\begin{omsources}
\item A.~Shleifer, ``Do demand curves for stocks slope down?'', \emph{Journal of Finance} 41(3), 1986.
\item L.~Harris and E.~Gurel, ``Price and volume effects associated with changes in the S\&P 500 list'', \emph{Journal of Finance} 41(4), 1986.
\item H.~Chen, G.~Noronha and V.~Singal, ``The price response to S\&P 500 index additions and deletions'', \emph{Journal of Finance} 59(4), 2004.
\item N.~Patel and I.~Welch, ``Extended stock returns in response to S\&P 500 index changes'', \emph{Review of Asset Pricing Studies} 7(2), 2017.
\item FTSE Russell (LSEG), Russell US indexes reconstitution announcements, 2026.
\end{omsources}

\section{Exercices}

\begin{exercise}[$\star$]\label{exo:s1:index-rebalancing:1}
Un indice de 300 actions a une bande de 15 rangs. À quel rang un non-membre entre-t-il dans l'indice, et à quel rang un membre en sort-il~?
\end{exercise}

\begin{exercise}[$\star$]\label{exo:s1:index-rebalancing:2}
Les \omterm{def:s1:index-rebalancing:fund}{fonds indiciels} doivent acheter 15~\% des actions d'une valeur~; celle-ci s'échange à hauteur de 0.6~\% de ses actions par jour, avec une volatilité quotidienne de 2~\%. Quelle poussée donne la loi de la racine carrée~?
\end{exercise}

\begin{exercise}[$\star$]\label{exo:s1:index-rebalancing:3}
Si la moitié de cette poussée est prise avant l'annonce et que l'autre moitié s'inverse après la date d'effet, que gagnent la transaction d'annonce et le fournisseur
de liquidité avant coûts~?
\end{exercise}

\begin{exercise}[$\star\star$]\label{exo:s1:index-rebalancing:4}
Pourquoi une bande réduit-elle la rotation, et que fait-elle à la prévisibilité des changements~?
\end{exercise}

\begin{exercise}[$\star\star$]\label{exo:s1:index-rebalancing:5}
Expliquer pourquoi le rendement d'événement lié à l'annonce diminue à mesure que le capital d'arbitrage croît, alors que l'événement lui-même ne se réduit pas.
\end{exercise}

\begin{exercise}[$\star\star$]\label{exo:s1:index-rebalancing:6}
Les prévisions faites soixante jours avant le \omterm{def:s1:index-rebalancing:fund}{jour de classement} sont justes dans 76~\% des cas. Que cela fait-il au rendement d'un prédicteur, par comparaison avec un prédicteur parfait~?
\end{exercise}

\begin{exercise}[$\star\star\star$]\label{exo:s1:index-rebalancing:7}
\emph{Programmation.} Exécuter \texttt{trades(0.30, 0.2)} et \texttt{trades(0.05, 0.8)}. Expliquer la différence dans le rendement de la transaction d'annonce.
\end{exercise}

\begin{exercise}[$\star\star\star$]\label{exo:s1:index-rebalancing:8}
\emph{Trouver la faille.} «~Les actions ajoutées à l'indice ont progressé de 7~\% entre l'annonce et la date d'effet dans notre échantillon des années 1990~; nous achèterons toutes les
additions annoncées.~»
\end{exercise}

\section{Problème~: Tout le monde le prévoit}

\begin{problem}[{Problème du week-end --- un événement connu, négocié}]\label{pb:s1:index-rebalancing:1}
L'indice synthétique et le registre public.

\medskip\noindent\textbf{Partie I --- Les règles.}
\begin{enumerate}
\item Définir un \omterm{def:s1:index-rebalancing:fund}{fonds indiciel}, un \omterm{def:s1:index-rebalancing:fund}{jour de classement} et une \omterm{def:s1:index-rebalancing:prediction}{prévision des entrées et sorties d'indice}.
\item Décrire le calendrier de Russell en 2026.
\item Décrire la règle et le calendrier de l'indice synthétique, ainsi que le nombre de changements.
\item Que fait une bande~?
\end{enumerate}

\medskip\noindent\textbf{Partie II --- Prévision.}
\begin{enumerate}[resume]
\item Comment les changements sont-ils prévus avant le \omterm{def:s1:index-rebalancing:fund}{jour de classement}~?
\item Donner la précision et le rappel à 0, 20 et 60 jours.
\item Pourquoi la prévision se dégrade-t-elle avec l'horizon~?
\item Quel est le compromis entre précoce et certain~?
\end{enumerate}

\medskip\noindent\textbf{Partie III --- L'événement.}
\begin{enumerate}[resume]
\item Décrire le chemin d'événement simulé.
\item Donner les rendements de la prévision, de la transaction d'annonce et de la fourniture de liquidité.
\item Que trouvent Shleifer, puis Chen, Noronha et Singal~?
\item Pourquoi les additions synthétiques conservent-elles une petite dérive~?
\end{enumerate}

\medskip\noindent\textbf{Partie IV --- Le verdict.}
\begin{enumerate}[resume]
\item Énoncer le \emph{résultat nommé}~: la précision de la prévision des changements de composition et le rendement d'événement lorsque la demande de réplication de l'indice croît par rapport au
capital d'arbitrage.
\item Que trouvent Patel et Welch~?
\item Où est passée la transaction~?
\item Quelle stratégie du fichier lancerait-on avec un petit capital, et laquelle avec un grand~?
\item Que faut-il surveiller à la date d'effet~?
\item Comment les reconstitutions semestrielles modifient-elles la transaction~?
\item Qu'est-ce qui distingue un événement d'indice d'un événement de flux du chapitre~9~?
\item En une phrase~: qu'est-ce que l'effet indice aujourd'hui~?
\end{enumerate}
\end{problem}

\section{Questions d'entretien}

\begin{interviewq}[$\star$ \iqroles{chercheur, trader}]\label{iq:s1:index-rebalancing:1}
Pourquoi le cours d'une action bouge-t-il lorsqu'elle entre dans un indice~?
\end{interviewq}

\begin{interviewq}[$\star\star$ \iqroles{trader}]\label{iq:s1:index-rebalancing:2}
Comment négocieriez-vous la reconstitution de Russell~?
\end{interviewq}

\begin{interviewq}[$\star\star$ \iqroles{chercheur, développeur}]\label{iq:s1:index-rebalancing:3}
Construisez un prédicteur des additions à un indice un mois avant le \omterm{def:s1:index-rebalancing:fund}{jour de classement}. Quelles entrées et quelle incertitude~?
\end{interviewq}

\begin{interviewq}[$\star\star$ \iqroles{risque}]\label{iq:s1:index-rebalancing:4}
Qu'est-ce qui peut mal tourner avec un portefeuille d'additions prévues à un indice~?
\end{interviewq}

\begin{interviewq}[$\star\star$ \iqroles{trader}]\label{iq:s1:index-rebalancing:5}
L'enchère de clôture le jour de la reconstitution montre un fort déséquilibre acheteur sur une action. Que faites-vous~?
\end{interviewq}

\begin{interviewq}[$\star\star\star$ \iqroles{chercheur}]\label{iq:s1:index-rebalancing:6}
La demande indicielle $D$ rencontre un capital d'arbitrage $A$ qui achète avant l'annonce. Dans un modèle où la poussée des cours est proportionnelle à la demande nette que le marché doit absorber
à chaque étape, comment le rendement après l'annonce dépend-il de $A/D$~?
\end{interviewq}
